% Generated draft: requires mathematical and source review.
\documentclass[11pt]{article}
\usepackage[margin=25mm]{geometry}
\usepackage{fontspec}
\setmainfont{Latin Modern Roman}
\usepackage{amsmath,amssymb,mathtools,unicode-math}
\setmathfont{Latin Modern Math}
\usepackage{enumitem,needspace,xurl,hyperref}
\setlength{\parindent}{0pt}
\setlength{\parskip}{5pt}
\setlength{\emergencystretch}{3em}
\newcommand{\N}{\mathbb{N}}
\newcommand{\Z}{\mathbb{Z}}
\newcommand{\Q}{\mathbb{Q}}
\newcommand{\R}{\mathbb{R}}
\newcommand{\C}{\mathbb{C}}
\newcommand{\F}{\mathbb{F}}
\newcommand{\A}{\mathbb{A}}
\newcommand{\PP}{\mathbb{P}}
\newcommand{\E}{\mathbb{E}}
\newcommand{\eps}{\varepsilon}
\newcommand{\dd}{\,\mathrm{d}}
\begin{document}
\section*{UnsolvedMath ID: 700773}
\textbf{Source problem code:} catalogue entry
\subsection{Definitions and mathematical statement}
Let \(G\) be a topological group with identity \(e\), and let \(\mathcal N_G(e)\) denote its neighborhoods of \(e\). For subsets \(A,B\subseteq G\), put \(AB=\{ab:a\in A,\ b\in B\}\). The group \(G\) is called \(F\)-balanced if
\[
\forall X\subseteq G\;\forall U\in\mathcal N_G(e)\;\exists V\in\mathcal N_G(e),\qquad VX\subseteq XU.
\]
Here \(V\) may depend on both \(X\) and \(U\). A subset \(W\subseteq G\) is invariant if \(gWg^{-1}=W\) for every \(g\in G\). The group \(G\) is called balanced if its identity has a neighborhood basis consisting of invariant sets; equivalently,
\[
\forall U\in\mathcal N_G(e)\;\exists W\in\mathcal N_G(e),\qquad W\subseteq U\quad\text{and}\quad \forall g\in G,\ gWg^{-1}=W.
\]
\textbf{Question (I. V. Protasov, 1990; Kourovka Notebook 11.81).} Is every \(F\)-balanced topological group balanced?

The supplied formulation does not specify a separation axiom for topological groups. Whether Hausdorffness is an implicit source convention remains unresolved.
\subsection{Short English statement}
If every subset \(X\) can absorb sufficiently small left multiplication into right multiplication by any prescribed identity neighborhood, must the group have a basis of conjugation-invariant identity neighborhoods? The source does not explicitly specify whether topological groups must be Hausdorff.
\subsection{Sources}
\begin{itemize}
\item Local UnsolvedMath registry, permanent ID \(700773\), exact source code \texttt{catalogue entry}, title \emph{Kourovka Notebook Problem 11.81}; supplied upstream revision \texttt{d371ee917ac9647fbb06cd6fc529deec6f3dda82}.
\item \emph{The Kourovka Notebook (Unsolved Problems in Group Theory)}, original problem 11.81, 11th issue (1990), proposed by I. V. Protasov. The supplied record identifies the 21st issue edition (2026), \texttt{arXiv:1401.0300v48}; its collection description names E. I. Khukhro and V. D. Mazurov as editors.
\item The supplied registry labels this problem \texttt{open} and records an update date of 2026-10-16. This is supplied registry metadata, not an independently verified assessment of the literature or mathematical status.
\end{itemize}
\end{document}
